\documentclass{article}
\usepackage{fontspec}
% switch to a monospace font supporting more Unicode characters
\setmonofont{JuliaMono}
\usepackage{listings}
\usepackage{amssymb, amsmath}

%colors for syntax highlighting
\usepackage{color}
\definecolor{keywordcolor}{rgb}{0.7, 0.1, 0.1}   % red
\definecolor{tacticcolor}{rgb}{0.0, 0.1, 0.6}    % blue
\definecolor{commentcolor}{rgb}{0.4, 0.4, 0.4}   % grey
\definecolor{symbolcolor}{rgb}{0.0, 0.1, 0.6}    % blue
\definecolor{sortcolor}{rgb}{0.1, 0.5, 0.1}      % green
\definecolor{attributecolor}{rgb}{0.7, 0.1, 0.1} % red

\def\lstlanguagefiles{lstlean.tex}
% set default language
\lstset{language=lean}

\begin{document}

\title{Report on the Workshop "Polyhedra in Lean"}
\author{
Georg Loho \\ Freie Universität Berlin \and
Olivia Röhrig \\ Zuse Institute Berlin \and
Martin Winter \\ Max Planck Institute Leipzig
}

\date{\today}

\maketitle

\section{Stub section}

Here's the hyperplane definition of a polytope over the reals, typeset using the \texttt{amsmath} package.

\[
  P = \left\{ x \in \mathbb{R}^n \;\middle|\; A x \le b \right\},
  \qquad A \in \mathbb{R}^{m \times n},\ b \in \mathbb{R}^m,
\]

Here's the point definition of a set being a polytope over a strictly ordered semiring in lean\cite{10.1007/978-3-030-79876-5_37}.

\begin{lstlisting}
variable [Semiring R] [PartialOrder R] [IsStrictOrderedRing R]
variable [ConvexSpace R X]

variable (R) in
/-- A set is a polytope if it is the convex hull of finitely many points. -/
def IsPolytope (s : Set X) : Prop := ∃ t : Finset X, s = convexHull R t
\end{lstlisting}

\bibliographystyle{plain}
\bibliography{references}

\end{document}
